\documentclass[14pt,aspectratio=169,xcolor=dvipsnames,table,onlytextwidth]{beamer}
\input{macros_en}
\csname endofdump\endcsname
\title{Bipartite Matching}
\subtitle{Introduction to Computational Mathematics}
\author{Instructor: \textbf{Tasuku Soma}}
\begin{document}
\maketitle

\begin{frame}<1>[label=outline]{Outline}
\begin{tikzpicture}
    \tikzstyle{block} = [rectangle, text width=\textwidth, rounded corners, minimum height=4em]
    \tikzstyle{line} = [draw, -latex]

    \node[block,alt={<1,3->{fill=cyan!10}{fill=red!10}}] (combopt) {
        \begin{columns}
            \begin{column}{0.5\linewidth}
        \textbf{1. Min-max theorem for bipartite matching}
            \end{column}
            \begin{column}{0.4\linewidth}
                \small
                \structure{\textbullet} Integrality of LP (recap)\\
                \structure{\textbullet} K\H{o}nig--Egerv\'ary theorem
                % \structure{\textbullet} Min-max theorem for weighted bipartite matching
            \end{column}
        \end{columns}
    };
    \node[block,alt={<1-2,4->{fill=cyan!10}{fill=red!10}}, below=.5em of combopt] (LP) {
        \begin{columns}
            \begin{column}{0.5\linewidth}
        \textbf{2. Unweighted bipartite matching}
            \end{column}
            \begin{column}{0.4\linewidth}
                \small
                \structure{\textbullet} Augmenting-path algorithm \\
                \structure{\textbullet} Revisiting the K\H{o}nig--Egerv\'ary theorem
            \end{column}
        \end{columns}
    };
    \node[block,alt={<1-3,5->{fill=cyan!10}{fill=red!10}}, below=.5em of LP] (TU) {
        \begin{columns}
            \begin{column}{0.5\linewidth}
        \textbf{3. Weighted bipartite matching}
            \end{column}
            \begin{column}{0.4\linewidth}
                \small
                \structure{\textbullet} Optimality condition \\
                \structure{\textbullet} Hungarian method
            \end{column}
        \end{columns}
    };
\end{tikzpicture}%
\end{frame}

%%%%%%%%%%%%%%%%%%%%%%%%%%%%%%%%%%%%%%%%%%%%%
\section{Min-max theorem for bipartite matching}
\againframe<2|handout:0>{outline}
%%%%%%%%%%%%%%%%%%%%%%%%%%%%%%%%%%%%%%%%%%%%%
\begin{frame}<1-3>
    \frametitle{Bipartite matching}

    \small
    \begin{columns}[T]
    \begin{column}{.6\textwidth}
        \begin{block}{Weighted bipartite matching problem}
            \structure{Input} $G=(V; E)$: a bipartite graph, \\ edge weights $w_e$ ($e \in E$) \\
            \structure{Output} A maximum-weight matching $M$ of $G$
        \end{block}


        \begin{block}<2->{\alt<2>{IP}{LP} formulation}
            \setlength{\abovedisplayskip}{0pt}
        \begin{alignat*}{2}
            \text{maximize} & \quad  \sum_{e \in E} w_e x_e \\
            \text{subject to} & \quad \sum_{e \in \delta(i)} x_e \leq 1 & \quad (i \in V) \\
            & \quad \alt<2>{x_e \in \{0,1\}}{} \only<3>{{\color{AlertOrange} x_e \geq 0}} & \quad (e \in E)
        \end{alignat*}
        \end{block}
        \only<2>{\footnotesize $\delta(i)$: the set of edges incident to vertex $i$}
    \end{column}
    \begin{column}{.4\textwidth}
        \centering
        \tikz{
            \graph[nodes={mynodes}, edges={myedges}, grow right=4em]{
            subgraph I_nm [V={1,2,3}, W={4,5}];
            1 -- {4,5};
            2 -- {4,5};
            3 -- {4};
            % matching
            {[edges={mypath}] {1,2} -- {4,5}; }
        };
        }
        \footnotesize
        \onslide<2->
        \begin{alignat*}{3}
            &\text{max}\quad  &&  w^\top x  \\
            &\text{s.t.}\quad  && x_{14} + x_{15}  &&\leq 1 \\
            & && x_{24} + x_{25}  &&\leq 1  \\
            & && x_{34}  &&\leq 1  \\
            & && x_{14} + x_{24} + x_{34} &&\leq 1  \\
            & && x_{15} + x_{25} &&\leq 1  \\
            & && 
            \only<3->{\color{AlertOrange}} x_{14}, x_{15}, x_{24}, x_{25}, x_{34} && \alt<2>{\in \{0, 1\} }{\color{AlertOrange}\geq 0}
        \end{alignat*}
    \end{column}
    \end{columns}
    \bigskip
    \onslide<3->{\structure{Recap:} Total unimodularity of the coefficient matrix ensures an LP optimal solution with \textbf{0--1} entries!}
\end{frame}

\begin{frame}
    \frametitle{Dual problem for bipartite matching}
    \small

    \begin{columns}[T]
    \begin{column}{.45\textwidth}
    \begin{primal}
        \setlength{\abovedisplayskip}{0pt}
        \begin{alignat*}{2}
            \text{max} & \quad  \sum_{e \in E} w_e x_e \\
            \text{s.t.} & \quad \sum_{e \in \delta(i)} x_e \leq 1  \quad (i \in V) \\
            & \quad x \geq 0
        \end{alignat*}
    \end{primal}
    \end{column}
    \begin{column}{.45\textwidth}
    \begin{dual}
        \setlength{\abovedisplayskip}{0pt}
        \begin{alignat*}{2}
            \text{min} & \quad  \sum_{i \in V} y_i \\
            \text{s.t.} & \quad y_i + y_j \geq w_e \quad (e=ij \in E) \\
            & \quad y \geq 0
        \end{alignat*}
    \end{dual}
    \end{column}
    \end{columns}

    \vfill
    The coefficient matrix of the dual problem (D) is also totally unimodular.
    Hence, if edge weights $w_e$ ($e \in E$) are integers, then (D) has an integral optimal solution.
\end{frame}

\begin{frame}
    \frametitle{K\H{o}nig--Egerv\'ary theorem}

    $G = (V^+, V^-; E)$: bipartite graph with vertex sets $V^+$ and $V^-$. 
    \begin{columns}
    \begin{column}{.75\textwidth}
    \begin{definition}
        $S \subseteq V^+, T \subseteq V^-$ is a \alert{vertex cover} \\
        $\defiff$ for any edge $e = ij \in E$, either $i \in S$ or $j \in T$.
    \end{definition}
    \end{column}
    \begin{column}{.25\textwidth}
        \centering
        \tikz{
            \graph[mygraph, grow right=4em]{
            subgraph I_nm [V={1,2,3}, W={4,5}];
            1 -- {4,5};
            2 -- {4,5};
            3 -- {4};
        };
        \begin{scope}[on background layer]
            \node[mysubset, fit=(1) (2), "$S$" left ]{} ;
            \node[mysubset, fit=(4), "$T$" right]{} ;
        \end{scope}
        \onslide<2->{
            \graph[edges={mypath}] {{(1),(2)} -- {(4), (5)}; };
        }
        }
        \footnotesize
    \end{column}
    \end{columns}

    \vfill
    \onslide<2->
    \begin{lemma}[Weak Duality]
        For any matching $M$ and vertex cover $(S,T)$, we have $\abs{M} \leq \abs{S} + \abs{T}$.
    \end{lemma}
    \begin{proof}
        Since $M$ consists of disjoint edges, at least $\abs{M}$ vertices are needed to cover $M$.
    \end{proof}
\end{frame}

\begin{frame}
    \frametitle{K\H{o}nig--Egerv\'ary theorem}
    Consider the previous LP with $w \equiv 1$. By integrality and strong duality, we obtain the following theorem.

    \bigskip
    \begin{theorem}[K\H{o}nig--Egerv\'ary]
        For any bipartite graph $G$, the following min-max relation holds:
        \[
            \max_{\text{$M$: matching}} |M| = \min_{\text{$(S,T)$: vertex cover}} |S| + |T|
        \]
    \end{theorem}
\end{frame}

\begin{frame}
    \frametitle{Example}

    \begin{columns}[T]
    \begin{column}{.7\textwidth}
        In the bipartite graph on the right,
        
        \begin{itemize}[<+(1)->]
            \item the maximum matching size is $|M| = 5$,
            \item the minimum vertex-cover size is $|S| + |T| = 5$.
        \end{itemize}

        \bigskip
        \onslide<+->{
        In general, to prove that a matching $M$ is maximum in a bipartite graph,
        it suffices to present a vertex cover $(S,T)$ satisfying
        \[ |M| = |S| + |T|. \]}
        \onslide<+->{
        Such a vertex cover is a certificate of optimality of $M$ \alert{(good characterization)}
        }
    \end{column}
    \begin{column}{.3\textwidth}
        \centering
        \tikz{
            \graph[mygraph, grow right=5em]{ mybipartite; };
        \onslide<2->{
            \graph[edges={mypath}] {{(1),(2),(3),(4),(5)} -- {(b),(c),(d),(e),(f)}; };
        }
        \begin{scope}[on background layer]
        \onslide<3->{
            \node[mysubset, fit=(1) (2), "$S$" left ]{} ;
            \node[mysubset, fit=(d) (e) (f), "$T$" right]{} ;
        }
        \end{scope}
        }
    \end{column}
    \end{columns}
\end{frame}

%%%%%%%%%%%%%%%%%%%%%%%%%%%%%%%%%%%%%%%%%%%%%
\section{Unweighted bipartite matching}
\againframe<3|handout:0>{outline}
%%%%%%%%%%%%%%%%%%%%%%%%%%%%%%%%%%%%%%%%%%%%%
\begin{frame}
    \frametitle{Augmenting paths}

    $M$: matching of $G$ \\
    \begin{definition}
        A path $P$ is an \alert{augmenting path} (with respect to $M$) $\defiff$
        \begin{enumerate}
            \item edges in $E \setminus M$ and in $M$ appear alternately in $P$.
            \item the two endpoints of $P$ are unmatched in $M$.
        \end{enumerate}
    \end{definition}
    \begin{center}
        \tikz[baseline=(1.base)]{\graph[mygraph]{
            1 -- 2 --[mypath, "{\footnotesize $\in M$}" {below, AlertOrange}] 3 --["{\footnotesize $\notin M$}" below] 4 --[mypath] 5 -- 6 --[mypath] 7 -- 8 ;
        }} \quad \textcolor{AlertOrange}{$\abs M = 3$}
        \pause
        \\ $\downarrow$ flip edges \\[1em]
        \tikz{\graph[mygraph]{
            1 --[mypath] 2 -- 3 --[mypath] 4 -- 5 --[mypath] 6 -- 7 --[mypath] 8 ;
        }} \quad \textcolor{AlertOrange}{$\abs M = 4$}
    \end{center}

    \vfill\pause
    Therefore, if $M$ is maximum, then no augmenting path exists.
\end{frame}

\begin{frame}
    \frametitle{Augmenting paths}

    \begin{lemma}
        If $M$ is not maximum, then an augmenting path exists.
    \end{lemma}

    \pause
    \newcommand{\M}{\textcolor{AlertOrange}{M}}
    \newcommand{\Mp}{\textcolor{UniGreen}{M'}}
    \begin{proof}
    Take two matchings $\M$, $\Mp$ with $\abs\M < \abs\Mp$.
    Then, each connected component of $\M + \Mp$ is either an alternating path or an even cycle.
    \begin{center}
        \tikz[]{\graph[mygraph]{1 --[bend left, mypath] 2; 1 --[bend right, mypath2] 2;}}
        \quad
        \tikz[baseline=(2.base)]{\graph[mygraph,clockwise]{subgraph I_n[n=4, radius=.7cm]; 
            {[edges={mypath }] 1 -- 2; 3 -- 4; };
            {[edges={mypath2}] 2 -- 3; 4 -- 1; };
        }}
        \quad
        \begin{tabular}{c}
        \tikz{\graph[mygraph]{
            1 --[mypath] 2 --[mypath2] 3 --[mypath] 4 --[mypath2] 5;
        }} \\
        \tikz{\graph[mygraph]{
            1 --[mypath] 2 --[mypath2] 3 --[mypath] 4;
        }} \\
        \tikz{\graph[mygraph]{
            1 --[mypath2] 2 --[mypath] 3 --[mypath2] 4 --[mypath] 5 --[mypath2] 6;
        }}
        \end{tabular}
    \end{center}
    Since $\abs\M < \abs\Mp$, some connected component of $\M + \Mp$ has more $\Mp$ edges than $\M$ edges. This is an augmenting path for $\M$. 
    \end{proof}
\end{frame}

\begin{frame}
    \frametitle{Augmenting-path algorithm}

    The lemma implies the following algorithm.

    \begin{block}{Augmenting-path algorithm}
        \begin{algorithmic}[1]
            \STATE Let $M \gets \emptyset$.
            \WHILE{there exists an augmenting path with respect to $M$}
                \STATE Find an augmenting path $P$.
                \STATE Flip edges along $P$ to obtain a larger matching $M$.
            \ENDWHILE
            \RETURN $M$
        \end{algorithmic}
    \end{block}

    \bigskip
    If an augmenting path can be found in $A$ time, then the whole algorithm runs in $O(nA)$ time.
    {\footnotesize \hfill ($n = |V|$)}
\end{frame}

\begin{frame}
    \frametitle{Finding augmenting paths via a directed graph}

    \small
    \begin{columns}[T]
    \begin{column}{.7\textwidth}
    For a matching $M$, orient edges of $G$ as follows:
    \begin{itemize}
        \item edges in $M$: both directions,
        \item edges in $E\setminus M$: from $V^+$ to $V^-$.
    \end{itemize}
    Let the resulting digraph be $D_M$.
    \onslide<2->{
        Also define
    \begin{itemize}
        \item $U^+ :=$ unmatched vertices in $V^+$ under $M$
        \item $U^- :=$ unmatched vertices in $V^-$ under $M$
    \end{itemize}
    }

    \onslide<3->{
    \begin{block}{}
        augmenting path $\longleftrightarrow$ directed path from $U^+$ to $U^-$ in $D_M$
    \end{block}
    Hence, BFS finds an augmenting path in $O(m)$ time. \\ \hfill {\footnotesize ($m = |E|$)}
    }
    \end{column}
    \begin{column}{.3\textwidth}
        \centering
        \tikz{
            \graph[mygraph, grow right=5em]{
            mybipartitedi;
            % matching
            {[edges={mypath}] {1,2,3,4} <-> {b,e,d,f}; }
        };
        \begin{scope}[on background layer]
        \onslide<2->{
            \node[mysubset, fit=(5) (6), "$U^+$" left ]{} ;
            \node[mysubset, fit=(a), "$U^-$" right]{} ;
            \node[mysubset, fit=(c), "$U^-$" right]{} ;
        }
        \end{scope}
        \node[above=10pt of 1]{$V^+$};
        \node[above=10pt of a]{$V^-$};
        }
        $D_M$
    \end{column}
    \end{columns}
\end{frame}

\begin{frame}
    \frametitle{Example}

    \begin{columns}[T]
    \begin{column}{.5\textwidth}
        \centering
        \vspace{5pt}

        \tikz{
        \graph[mygraph, grow right=5em]{mybipartite};
            \onslide<3|handout:0>{\graph[edges={mypath}]{
                (1) -- (b);
            };}
            \onslide<4|handout:0>{\graph[edges={mypath}]{
                {(1),(2)} -- {(a),(b)};
            };}
            \onslide<5|handout:0>{\graph[edges={mypath}]{
                {(1),(2),(3),(4)} -- {(a),(b),(f),(e)};
            };}
            \onslide<6>{\graph[edges={mypath}]{
                {(1),(2),(3),(4),(5)} -- {(a),(b),(d),(e),(f)};
            };}
        } \\ $G$
    \end{column}
    \begin{column}{.5\textwidth}
        \centering
        \onslide<2->{
        \tikz{
            \graph[mygraph, grow right=5em]{mybipartitedi};
            \onslide<3|handout:0>{\graph[edges={mypath}]{
                (1) <-> (b);
            };}
            \onslide<4|handout:0>{\graph[edges={mypath}]{
                {(1),(2)} <-> {(a),(b)};
            };}
            \onslide<5|handout:0>{\graph[edges={mypath}]{
                {(1),(2),(3),(4)} <-> {(a),(b),(f),(e)};
            };}
            \onslide<6>{\graph[edges={mypath}]{
                {(1),(2),(3),(4),(5)} <-> {(a),(b),(d),(e),(f)};
            };}
        \begin{scope}[on background layer]
        \onslide<2|handout:0>{
            \node[mysubset, fit=(1) (2) (3) (4) (5) (6), "$U^+$" left ]{} ;
            \node[mysubset, fit=(a) (b) (c) (d) (e) (f), "$U^-$" right]{} ;
        }
        \onslide<3|handout:0>{
            \node[mysubset, fit=(2) (3) (4) (5) (6), "$U^+$" left ]{} ;
            \node[mysubset, fit=(a), "$U^-$" right]{} ;
            \node[mysubset, fit=(c) (d) (e) (f), "$U^-$" right]{} ;
        }
        \onslide<4|handout:0>{
            \node[mysubset, fit=(3) (4) (5) (6), "$U^+$" left ]{} ;
            \node[mysubset, fit=(c) (d) (e) (f), "$U^-$" right]{} ;
        }
        \onslide<5|handout:0>{
            \node[mysubset, fit=(5) (6), "$U^+$" left ]{} ;
            \node[mysubset, fit=(c) (d), "$U^-$" right]{} ;
        }
        \onslide<6>{
            \node[mysubset, fit=(6), "$U^+$" left ]{} ;
            \node[mysubset, fit=(c), "$U^-$" right]{} ;
        }
        \end{scope}
        } \\ $D_M$
        }
    \end{column}
    \end{columns}
    \hfill \onslide<6>{[Finish]}
\end{frame}

\begin{frame}
    \frametitle{Revisiting K\H{o}nig--Egerv\'ary}
    \postit{
        \footnotesize
    $\displaystyle
            \max_{\text{$M$:matching}} |M| = \min_{\text{$(S,T)$: vertex cover}} |S| + |T|
            $
    }
    \small
    \begin{columns}[T]
    \begin{column}{.7\textwidth}
    \vspace{-2em}
    \begin{theorem}
        At the termination of the algorithm, let $X$ be the set of all vertices reachable from $U^+$ in $D_M$. Define
        \setlength{\abovedisplayskip}{5pt}
        \setlength{\belowdisplayskip}{5pt}
        \[
            S := V^+ \setminus X, \quad T := V^- \cap X.
        \]
        Then $(S,T)$ is a minimum vertex cover.
    \end{theorem}

    \onslide<2->{
    \begin{proof} 
        \structure{Why vertex cover?} One can traverse any edge from left to right in $D_M$, so $(S,T)$ is a vertex cover by definition.

        \medskip
        \structure{Why minimum?}
        Matching edges are bidirected, so their endpoints are either both in $X$ or both not in $X$.
        Since no augmenting path exists, $X \cap U^- = \emptyset$.
        Hence every vertex in $(S,T)$ covers exactly one edge of $M$. Thus, $|S|+|T|=|M|$.
    \end{proof}
    }
    \end{column}
    \begin{column}{.3\textwidth}
        \centering
        \tikz{
            \graph[mygraph, grow right=5em]{
            mybipartitedi;

            % matching
            {[edges={mypath}] {1,2,3,4,5} <-> {a,b,d,e,f}; }
        };
        \begin{scope}[on background layer]
            \path[draw,AlertOrange!10,fill=AlertOrange!10,line width=1.0cm, rounded corners] (5.north west) -- (6.south west) -- (f.south west) -- (f.north west) -- cycle node[AlertOrange,left]{$X$};
            \node[mysubset, fit=(6), "$U^+$" left ]{} ;
            \node[mysubset, fit=(c), "$U^-$" right]{} ;
            \node[draw,UniPink,dotted,thick,fill=UniPink!10, fit=(1) (2) (3) (4), "$S$" left ]{} ;
            \node[draw,UniPink,dotted,thick,fill=UniPink!10, fit=(f), "$T$" right]{} ;
        \end{scope}
        }
    \end{column}
    \end{columns}

    \centering
    \bigskip
    \onslide<3->{\bf An algorithmic proof of the K\H{o}nig--Egerv\'ary theorem without LP!}
\end{frame}

\begin{frame}
    \frametitle{Unweighted bipartite matching (summary)}

    \begin{theorem}
        In a bipartite graph:
        \begin{itemize}
            \setlength{\itemsep}{1em}
            \item The augmenting-path algorithm finds a maximum matching $M$ in $O(mn)$ time. \hfill {\footnotesize ($m = |E|$, $n = |V|$)}

            \item It also finds a vertex cover $(S,T)$ satisfying $|M| = |S| + |T|$.
        \end{itemize}
    \end{theorem}
\end{frame}

%%%%%%%%%%%%%%%%%%%%%%%%%%%%%%%%%%%%%%%%%%%%%
\section{Weighted bipartite matching}
\againframe<4|handout:0>{outline}
%%%%%%%%%%%%%%%%%%%%%%%%%%%%%%%%%%%%%%%%%%%%%
\begin{frame}
    \frametitle{Weighted bipartite (perfect) matching}
    \small

    A matching $M$ is \alert{perfect} $\defiff$ every vertex is incident to $M$ $\iff$ $2|M| = |V|$.

    \begin{block}{Weighted perfect bipartite matching problem}
        \structure{Input} $G=(V; E)$: bipartite graph ($|V^+|=|V^-|=n/2$), edge weights $w_e$ ($e \in E$) \\
        \structure{Output} A perfect matching $M$ maximizing $w(M) := \sum_{e\in M}w_e$
    \end{block}
    Assume that $G$ has at least one perfect matching in what follows.
    \pause

    \begin{columns}[T]
    \begin{column}{.45\textwidth}
    \begin{primal}
        \setlength{\abovedisplayskip}{0pt}
        \begin{alignat*}{2}
            \text{max} & \quad  \sum_{e \in E} w_e x_e \\
            \text{s.t.} & \quad \sum_{e \in \delta(i)} x_e = 1  \quad (i \in V) \\
            & \quad x \geq 0
        \end{alignat*}
    \end{primal}
    \end{column}
    \begin{column}{.45\textwidth}
    \begin{dual}
        \setlength{\abovedisplayskip}{0pt}
        \begin{alignat*}{2}
            \text{min} & \quad  \sum_{i \in V} y_i =: y(V) \\
            \text{s.t.} & \quad y_i + y_j \geq w_e \quad (e=ij \in E)
        \end{alignat*}
    \end{dual}
    \end{column}
    \end{columns}

\end{frame}

\begin{frame}
    \frametitle{Optimality Condition}

    \begin{columns}[T]
    \begin{column}{.6\textwidth}
        \begin{lemma}[Optimality condition]
        $M \subseteq E$, $y \in \R^V$ are optimal solutions
        $\iff$
        \begin{enumerate}
            \item $M$ is a perfect matching,
            \item $y$ is feasible for (D),
            \item $y_i + y_j = w_e$ \quad ($e = ij \in M$)
        \end{enumerate}
        \end{lemma}
        
        \onslide<2->{
                \begin{itemize}
                    \item a feasible solution of (D) is called a \alert{$w$-cover},
                    \item edges satisfying condition (3) are said to be \alert{tight} (w.r.t.~$y$).
                \end{itemize}
        }
    \end{column}
    \begin{column}{.35\textwidth}
        \footnotesize
        \begin{primal}
        \begin{alignat*}{2}
            \text{max} & \quad  \sum_{e \in E} w_e x_e \\
            \text{s.t.} & \quad \sum_{e \in \delta(i)} x_e = 1  \quad (i \in V) \\
            & \quad x \geq 0
        \end{alignat*}
    \end{primal}
    \begin{dual}
        \setlength{\abovedisplayskip}{0pt}
        \begin{alignat*}{2}
            \text{min} & \quad  \sum_{i \in V} y_i =: y(V) \\
            \text{s.t.} & \quad y_i + y_j \geq w_e \quad (e \in E)
        \end{alignat*}
    \end{dual}
    \end{column}
    \end{columns}

\end{frame}

\begin{frame}
    \frametitle{Hungarian method}

    \small
    The algorithm maintains a matching $M$ and a $w$-cover $y$ satisfying (2)(3), and updates them until (1) is satisfied.

    \begin{definition}
        For a $w$-cover $y$, define a subgraph $G_y = (V^+, V^-, E_y)$ of $G$ by
        \[
            E_y = \{e = ij \in E : y_i + y_j = w_e \}
        \]
        That is, $G_y$ keeps only tight edges with respect to $y$.
    \end{definition}

    \begin{lemma}
        If $G_y$ has a \textbf{perfect} matching, then it is a maximum-weight matching.
    \end{lemma}
    \begin{proof}
        Optimality condition!
    \end{proof}

\end{frame}

\begin{frame}
    \frametitle{Hungarian method}

    \small
    \begin{columns}[T]
    \begin{column}{.7\textwidth}
    What if $G_y$ has no perfect matching?\\
    $\longrightarrow$ We can update $y$ and improve the objective value of (D).
    
    \pause
    \begin{lemma}
        Suppose $G_y$ has no perfect matching.
        Let $(S, T)$ be a minimum vertex cover of $G_y$.
        Define
        $\eps := \min\{ y_i + y_j - w_e : e = ij \in E, i \notin S, j \notin T \}$.
        Let $y'$ be
            \setlength{\abovedisplayskip}{0pt}
        \[
        y'_i =
        \begin{cases}
            y_i - \eps & (i \in V^+ \setminus S) \\
            y_i + \eps & (i \in V^- \cap T) \\
            y_i        & (\text{otherwise})
        \end{cases}
        \]
        Then $y'$ is a $w$-cover with $y'(V) < y(V)$.
    \end{lemma}
    \end{column}
    \begin{column}{.3\textwidth}
    \begin{center}
    \tikz{\graph[mygraph,simple,grow right=5em, branch down=3em]{
        {[edges=gray!20] mybipartite2;};
    1 -> b, 2 -> {a,b}, 3 -> a;
    {[edges=mypath] {(1),(2)} <-> {(b),(a)}};
    };
    \foreach \y\i in {4/1,2/2,3/3}
        \node[UniGreen,left=1pt of \i]{$\y$};
    \foreach \y\i in {2/a,0/b,0/c}
        \node[UniGreen,right=1pt of \i]{$\y$};
    \begin{scope}[on background layer]
        % \path[draw,AlertOrange!10,fill=AlertOrange!10,line width=1.0cm, rounded corners, visible on=<.->] (1.north west) -- (3.south west) -- (b.south east) -- (a.north east) -- cycle node[AlertOrange,left]{$X$};
        \node[draw,UniPink,dotted,thick,fill=UniPink!10, fit=(a) (b), "$T$" right]{} ;
    \end{scope}
        \foreach \i in {1,2,3}
            \node[UniGreen,above] at (\i) {$-\eps$};
        \foreach \i in {a,b}
            \node[UniGreen,above] at (\i) {$+\eps$};
    } 
    \end{center}
    \hspace{1em} $S = \emptyset$
    \end{column}
    \end{columns}

\end{frame}

\begin{frame}
    \frametitle{Proof}
    \postit{
        \footnotesize
        $\eps := \min\{ y_i + y_j - w_e : e = ij \in E, i \notin S, j \notin T \}$.
        \\
        \footnotesize
            $
        y'_i =
        \begin{cases}
            y_i - \eps & (i \in V^+ \setminus S) \\
            y_i + \eps & (i \in V^- \cap T) \\
            y_i        & (\text{otherwise})
        \end{cases}
        $
    }

    \small

    \structure{Why $y'$ is a $w$-cover?}
    \begin{itemize}
        \item The only constraints $y_i + y_j \geq w_e$ ($e = ij \in E$) that could become violated after the update are those with $i \notin S$, $j \notin T$.
        \item By the definition of $\eps$, feasibility is preserved.
    \end{itemize}

    \bigskip\pause
    \structure{Why $y'(V) < y(V)$?}
    \begin{itemize}
        \item $y'(V) = y(V) - \eps(|V^+ \setminus S| - |T|) = y(V) - \eps(n/2 - |S| - |T|)$.

        \item Since $G_y$ has no perfect matching, we have $|S| + |T| < n/2$ by K\H{o}nig--Egerv\'ary theorem.

        \item Thus it suffices to show $\eps > 0$. Because $(S,T)$ is a vertex cover of $G_y$, any edge $e = ij \in E$ with $i \notin S$, $j \notin T$ is not tight, i.e., $y_i + y_j - w_e > 0$. \qed
    \end{itemize}
\end{frame}

\begin{frame}
    \frametitle{Hungarian method}

    \small
    \begin{block}{}
        \begin{algorithmic}[1]
            \STATE Set $y_i := \max_{e \in \delta(i)} w_e$ ($i \in V^+$), $y_j := 0$ ($j \in V^-$).
            \COMMENT{$y$ is a trivial $w$-cover}
            \WHILE{\textbf{True}}
                \STATE Find a maximum (unweighted) matching $M$ of $G_y$.
                \COMMENT{$O(mn)$ time}
                \IF{$M$ is perfect}
                    \RETURN $M$
                    \COMMENT{By optimality condition, $M$ is maximum-weight}
                \ELSE[{$w$-cover update}]
                    \STATE Let $X$ be the set of vertices reachable from $U^+$ in $D_M(y)$. \\
                    \COMMENT{$(S, T) = (V^+ \setminus X, V^- \cap X)$ is a minimum vertex cover of $G_y$}
                    \STATE $\eps := \min\{ y_i + y_j - w_{ij} : i \in V^+ \cap X, j \in V^- \setminus X \}$
                    \STATE $y_i \gets y_i - \eps$ ($i \in V^+ \cap X$), $y_j \gets y_j + \eps$ ($j \in V^- \cap X$)
                \ENDIF
            \ENDWHILE
        \end{algorithmic}
    \end{block}
\end{frame}

\begin{frame}
    \frametitle{Example}
    \small
    \begin{columns}
    \begin{column}{.20\textwidth}<+->
        \centering
        \tikz{\graph[mygraph,grow right=5em, branch down=3em]{mybipartite2;
            1 --["{$5$}" above] a;
            1 --["{$4$}" {above, pos=.3}] b;
            2 --["{$4$}" {above, pos=.1}] a;
            2 --["{$2$}" {below, pos=.9}] b;
            2 --["{$1$}" {above, pos=.99}] c;
            3 --["{$5$}" {above, pos=.02}] a;
            3 --["{$2$}" below] c;
        };}
        \\ $G$
    \end{column}
    \begin{column}{.25\textwidth}<+->
        \centering
        \tikz{\graph[mygraph,simple,grow right=5em, branch down=3em]{
        {[edges=gray!20]
            mybipartite2;
        };
        {1,2,3} -> a;
        };
        \foreach \y\i in {5/1,4/2,5/3}
            \node[UniGreen,left=1pt of \i]{$\y$};
        \foreach \y\i in {0/a,0/b,0/c}
            \node[UniGreen,right=1pt of \i]{$\y$};
        \onslide<+->{\graph[edges=mypath]{(1) <-> (a)};}
        \begin{scope}[on background layer]
            \path[draw,AlertOrange!10,fill=AlertOrange!10,line width=1.0cm, rounded corners,visible on=<+->] (1.north west) -- (3.south west) -- (a.south east) -- (a.north east) -- cycle node[AlertOrange,left]{$X$};
        \end{scope}
        \onslide<+->{
            \foreach \i in {1,2,3}
                \node[UniGreen,above] at (\i) {$-\eps$};
            \node[UniGreen,above] at (a) {$+\eps$};
        }
        } $y(V)=14$, \onslide<.->{$\eps=1$}
    \end{column}
    \begin{column}{.25\textwidth}<+->
        \centering
        \tikz{\graph[mygraph,simple,grow right=5em, branch down=3em]{
            {[edges=gray!20] mybipartite2;};
        1 -> {a,b}, 2 -> a; 3 -> a;
        };
        \onslide<.>{
            \graph[edges=mypath]{(1) <-> (a);};
        }
        \foreach \y\i in {4/1,3/2,4/3}
            \node[UniGreen,left=1pt of \i]{$\y$};
        \foreach \y\i in {1/a,0/b,0/c}
            \node[UniGreen,right=1pt of \i]{$\y$};
        \onslide<+->{
        \graph[edges=mypath]{{(1),(2)} <-> {(b),(a)}};
        \begin{scope}[on background layer]
            \path[draw,AlertOrange!10,fill=AlertOrange!10,line width=1.0cm, rounded corners,visible on=<.->] (2.north west) -- (3.south west) -- (a.south east) -- (a.north east) -- cycle node[AlertOrange,left]{$X$};
        \end{scope}
            \foreach \i in {2,3}
                \node[UniGreen,above] at (\i) {$-\eps$};
            \node[UniGreen,above] at (a) {$+\eps$};
        }
        } $y(V)=12$, \onslide<.->{$\eps=1$}
    \end{column}
    \begin{column}{.25\textwidth}<+->
        \centering
        \tikz{\graph[mygraph,simple,grow right=5em, branch down=3em]{
            {[edges=gray!20] mybipartite2;};
        1 -> b, 2 -> {a,b}, 3 -> a;
        {[edges=mypath] {(1),(2)} <-> {(b),(a)}};
        };
        \foreach \y\i in {4/1,2/2,3/3}
            \node[UniGreen,left=1pt of \i]{$\y$};
        \foreach \y\i in {2/a,0/b,0/c}
            \node[UniGreen,right=1pt of \i]{$\y$};
        \onslide<+->{
        \begin{scope}[on background layer]
            \path[draw,AlertOrange!10,fill=AlertOrange!10,line width=1.0cm, rounded corners, visible on=<.->] (1.north west) -- (3.south west) -- (b.south east) -- (a.north east) -- cycle node[AlertOrange,left]{$X$};
        \end{scope}
            \foreach \i in {1,2,3}
                \node[UniGreen,above] at (\i) {$-\eps$};
            \foreach \i in {a,b}
                \node[UniGreen,above] at (\i) {$+\eps$};
        }
        } $y(V) = 11$, \onslide<.->{$\eps=1$}
    \end{column}
    \end{columns}

    \begin{columns}
    \begin{column}{.35\textwidth}<+->
        \tikz{\graph[mygraph,simple,grow right=5em, branch down=3em]{
            {[edges=gray!20] mybipartite2;};
        1 -> b, 2 -> {a,b,c}, 3 -> {a,c};
        {[edges=mypath] {1,2} <-> {b,a};};
        };
        \foreach \y\i in {3/1,1/2,2/3}
            \node[UniGreen,left=1pt of \i]{$\y$};
        \foreach \y\i in {3/a,1/b,0/c}
            \node[UniGreen,right=1pt of \i]{$\y$};
        \onslide<+->{
        \graph[edges=mypath]{{(1),(2),(3)} <-> {(b),(a),(c)}};
        }
        } $y(V) = 10$\onslide<.->{ $=w(M)$ [END]}
    \end{column}
    \begin{column}{.65\textwidth}

    \end{column}
    \end{columns}

\end{frame}

\begin{frame}{Analysis of the Hungarian method}

    \small
    \begin{theorem}
        The Hungarian method finds a maximum-weight perfect matching and a minimum $w$-cover in $O(mn^3)$ time.
        \hfill {\footnotesize ($m = |E|$, $n = |V|$)}
    \end{theorem}

    \pause
    \begin{itemize}[<+->]
        \item During the algorithm, $|M|$ never decreases.
        \begin{itemize}
            \item[($\because$)] when $y$ is updated to $y'$, the maximum matching $M$ in $G_y$ is still a matching in $G_{y'}$.
        \end{itemize}
        $\longrightarrow$ $M$ is updated at most $n/2$ times.

        \item When $y$ is updated without increasing $|M|$, the edges achieving the minimum in the definition of $\eps$ make the reachable set $X$ strictly larger.\\
        $\longrightarrow$ after at most $n$ such updates of $y$, $|M|$ must increase.
    \end{itemize}

    \vfill
    \onslide<+->
    Therefore, the total number of iterations is $O(n^2)$. \qed
\end{frame}

\begin{frame}
    \frametitle{Summary}

    \begin{itemize}
        \setlength{\itemsep}{1em}
        \item Basic LP-based ideas in combinatorial optimization (integral polyhedra and totally unimodular matrices)

        \item Using LP integrality, we proved the min-max theorem for bipartite matching (K\H{o}nig--Egerv\'ary theorem)

        \item Combinatorial algorithms that do not use LP
        \begin{itemize}
            \item augmenting-path algorithm for unweighted bipartite matching
            \item Hungarian method for weighted bipartite perfect matching
        \end{itemize}
    \end{itemize}

\end{frame}

\end{document}
